\chapter{Engenharia de Software e Modernização}
\label{sec:engenhariasoftware}

\section{Engenharia de Requisitos com Sommerville}

Sommerville~\cite{sommerville2011} estrutura a engenharia de requisitos em quatro fases interdependentes: elicitação, análise, especificação e validação. A reestruturação do sistema em três módulos — \texttt{models.py}, \texttt{components.py} e \texttt{app.py} — evidencia como cada fase foi aplicada para preservar as regras de negócio do legado Vaadin/Java sem reaproveitar seu código-fonte.

\subsection{Elicitação}
As entradas críticas do domínio de regressão espectral foram identificadas diretamente no sistema Java original: as matrizes \texttt{Xcal} e \texttt{Xteste}, os vetores \texttt{Ycal} e \texttt{Yteste}, e os hiperparâmetros \textit{m} (número de \textit{bags}) e \textit{k} (fração de \textit{subagging}). O \texttt{app.py} captura esses insumos por meio de \texttt{st.file\_uploader} e \texttt{st.text\_input}, incluindo suporte a decimal com vírgula — particularidade do padrão brasileiro — tratada nas funções \texttt{parse\_matrix()} e \texttt{parse\_vector()} do \texttt{models.py}.

\subsection{Análise}
A análise revelou que a lógica do \texttt{RegressionService.java} podia ser decomposta em responsabilidades distintas: \textit{parse} de arquivos, algoritmos de regressão, métricas de avaliação e otimização por validação cruzada. Essa decomposição originou a estrutura modular adotada, onde \texttt{models.py} concentra toda a inteligência matemática, sem qualquer dependência de interface.

\subsection{Especificação}
As assinaturas de entrada e saída das principais funções do~\verb|models.py|~foram documentadas com \textit{type hints} \footnote{\textit{Type hints} são anotações opcionais da linguagem Python que declaram os tipos esperados dos parâmetros e do valor de retorno de uma função
(por exemplo, \texttt{x: np.ndarray} e \texttt{-> float}). Documentam a assinatura da função e permitem verificação estática por ferramentas externas, mas não são impostas pelo interpretador em tempo de execução.} nas principais funções do \texttt{models.py}:
\texttt{parse\_matrix(file\_obj) -> np.ndarray} e
\texttt{parse\_vector(file\_obj) -> np.ndarray}, em que \texttt{file\_obj} é um
objeto de arquivo (\textit{file-like}) que expõe \texttt{.name} e \texttt{.read()},
sendo a extensão (\texttt{.txt}, \texttt{.csv} ou \texttt{.xlsx}) usada para
escolher o \textit{parser};
\texttt{perform\_rlm(Xcal, Ycal, Xteste: np.ndarray) -> np.ndarray};
\texttt{calc\_metrics(actual, predicted) -> tuple[float, float, float]}.
A função \texttt{run\_full\_analysis()} retorna um dicionário (\texttt{dict})
com as chaves \texttt{log}, \texttt{optimal}, \texttt{metrics} e \texttt{preds},
definindo a interface de comunicação entre a camada de negócio e a camada de apresentação.

\subsection{Validação}
A validação cruzada 5-\textit{fold} implementada em \texttt{run\_cross\_validation()} opera como mecanismo de validação de equivalência, permitindo comparar o desempenho do sistema Python com os resultados históricos do sistema Java. O \textit{log} estruturado gerado por \texttt{run\_full\_analysis()} documenta cada etapa — dimensões detectadas, $\text{RMSEP}_{\text{CV}}$ por combinação de hiperparâmetros, métricas finais e tempo de execução — possibilitando auditoria científica e rastreabilidade completa. Para assegurar a estabilidade dessa arquitetura, construiu-se uma suíte com 92 testes automatizados, composta por 21 testes unitários, 60 testes de caixa-preta e 11 testes de carga e estresse.

\section{Arquitetura Modular em Python/Streamlit}

A arquitetura do sistema modernizado adota o padrão de separação de responsabilidades em três camadas físicas distintas, representadas pelos arquivos \texttt{models.py}, \texttt{components.py} e \texttt{app.py}. Essa organização substitui a estrutura monolítica do sistema Java original, garantindo testabilidade isolada e desacoplamento entre a lógica de negócio e a interface com o usuário.

O módulo \texttt{models.py} atua como a lógica matemática do sistema, encapsulando todas as operações matemáticas e estatísticas sem qualquer importação de bibliotecas de interface gráfica. Ele contém as implementações da RLM via \texttt{np.linalg.lstsq}, do \textit{bagging} e do \textit{subagging} parametrizados, das funções de \textit{parsing} de dados e das rotinas de validação cruzada 5-\textit{fold} com \textit{grid search}.

Ressalta-se que o sistema Java legado calcula os coeficientes da RLM pela
pseudoinversa de Moore-Penrose explícita (\texttt{SimpleMatrix.pseudoInverse()},
biblioteca EJML, implementada em Java puro), ao passo que a versão Python emprega
\texttt{np.linalg.lstsq}, que delega ao NumPy a rotina \texttt{gelsd} da
biblioteca LAPACK~\cite{anderson1999lapack}. Essa rotina resolve o problema de mínimos quadrados por
decomposição em valores singulares (SVD) e retorna a solução de mínima norma.
Para os sistemas subdeterminados deste trabalho ($17 \times 723$), ambas as
rotinas convergem para a mesma solução teórica
$\widehat{\boldsymbol{\beta}} = \mathbf{X}^{+}\mathbf{y}$, embora por caminhos
numéricos distintos.

O módulo \texttt{components.py} responde pela fábrica de componentes visuais e geração de gráficos interativos utilizando a biblioteca Plotly. Ele transforma as matrizes de resultados e vetores de predição retornados pela camada de modelagem em objetos de figura prontos para renderização, isolando o código de visualização do fluxo principal da aplicação.

O módulo \texttt{app.py} funciona como controlador da interface e do estado da sessão (\texttt{st.session\_state}). Ele captura as interações do usuário, invoca as funções puras de \texttt{models.py}, repassa os resultados para \texttt{components.py} e gerencia a exibição sequencial em abas. Essa divisão garante que alterações na camada visual não impactem os algoritmos científicos e vice-versa.

%A Figura~\ref{fig:arquitetura_python} apresenta as três camadas do sistema modernizado, em contraste com a organização monolítica do sistema Java original, ilustrada na Figura~\ref{fig:arquitetura_java}. O diagrama de classes dos módulos (Figura~\ref{fig:BaggingOptimized_classe3}) e os demais diagramas \gls{UML} estão reunidos no Apêndice~\ref{cha:testesgraficos}.

A Figura~\ref{fig:arquitetura_java} mostra a organização lógica em camadas do sistema Java original, contida em um único arquivo, e a Figura~\ref{fig:arquitetura_python} apresenta as três camadas físicas do sistema modernizado. O diagrama de classes dos módulos (Figura~\ref{fig:BaggingOptimized_classe3}) e os demais diagramas \gls{UML} estão no Apêndice~\ref{cha:testesgraficos}.

\begin{figure}[H]
    \caption{Organização lógica em camadas dentro de um único arquivo --- Sistema legado monolítico (Java Desktop).}
    \centering
    \begin{tikzpicture}[
        layer/.style={draw, rounded corners, minimum width=8cm,
                      minimum height=1.2cm, align=center, fill=#1},
        arrow/.style={-{Stealth}, thick}
    ]
        \node[layer=blue!15] (view) {Camada de Apresentação\\\texttt{Swing/AWT} (\texttt{JFrame}, \texttt{JPanel}, \textit{listeners})};
        \node[layer=green!15, below=0.8cm of view] (controller) {Camada de Controle\\Orquestração de eventos (\texttt{ActionListener})};
        \node[layer=orange!15, below=0.8cm of controller] (model) {lógica matemática\\\texttt{BaggingOptimized.java} + EJML (\texttt{SimpleMatrix})};

        \draw[arrow] (view) -- (controller);
        \draw[arrow] (controller) -- (model);
    \end{tikzpicture}
    \caption*{Fonte: \oautor}
    \label{fig:arquitetura_java}
\end{figure}

\begin{figure}[H]
    \caption{Arquitetura em camadas --- Sistema modernizado (Python/Streamlit).}
    \centering
    \begin{tikzpicture}[
        layer/.style={draw, rounded corners, minimum width=8cm,
                      minimum height=1.2cm, align=center, fill=#1},
        arrow/.style={-{Stealth}, thick}
    ]
        \node[layer=blue!15] (apresentacao) {Camada de Apresentação\\\texttt{app.py} --- Streamlit (orquestração de UI e sessão)};
        \node[layer=green!15, below=0.8cm of apresentacao] (visual) {Camada de Visualização\\\texttt{components.py} --- Plotly (fábrica de gráficos)};
        \node[layer=orange!15, below=0.8cm of visual] (nucleo) {lógica matemática\\\texttt{models.py} --- NumPy (\texttt{np.linalg.lstsq})};

        \draw[arrow] (apresentacao) -- (visual);
        \draw[arrow] (visual) -- (nucleo);
    \end{tikzpicture}
    \caption*{Fonte: \oautor}
    \label{fig:arquitetura_python}
\end{figure}

